% ============================================================================
\section{Introduction}\label{sec:intro}
% ============================================================================

Fresh food is replenished daily, its stock ages every night, and whatever is not sold within its shelf life is thrown away. Retailers and their distributors therefore face the classical perishable inventory trade-off between waste and lost sales \citep{nahmias1982perishable,bakker2012review}, but with two complications that single-store inventory models do not capture. First, deliveries to different stores share a small fleet: a store that is not on today's route cannot be replenished today, and a vehicle that is full cannot top up the next store. Second, retailers now have rich contextual information---promotion calendars, weather forecasts, recent sales---that machine-learning forecasters turn into store-day demand predictions. The operational question is how such forecasts should be translated into replenishment and delivery decisions when stock ages, shelf space is limited and vehicles are shared.

Two literatures answer different halves of this question. Perishable inventory control offers simple, robust ordering rules---base-stock policies \citep{bu2023asymptotic}, and rules that anticipate the waste of the current stock \citep{broekmeulen2009heuristic,haijema2019improved}---but treats each store in isolation. Inventory routing integrates replenishment and routing, with exact methods for deterministic perishable demand \citep{coelho2014optimal,alvarez2020formulations}, decomposition and order-up-to rules for stochastic perishable demand \citep{crama2018stochastic}, robust formulations \citep{solyali2012robust,he2024robust} and, recently, scenario-based predict-then-optimise policies that use machine-learning quantile forecasts \citep{jia2025scenario}. In parallel, decision-focused learning (DFL) trains forecasters on the cost of the decisions they induce rather than on forecast accuracy \citep{elmachtoub2022smart,mandi2024decision}. Whether DFL adds anything once the strong inventory rules and stochastic lookahead methods are on the table has not been tested for perishable inventory routing.

This paper studies a periodic-review perishable inventory-routing system with contextual demand, lost sales, FIFO issuing and shelf capacity, and compares twelve replenishment policies in closed loop, every day observing the age-structured stock and the same forecasts. The comparison includes the inventory-control rules a perishable-inventory manager would use (base-stock with and without estimated waste), the deliver-up-to-level rule of \citet{crama2018stochastic}, and rolling-horizon inventory-routing models driven by mean, quantile, robust and scenario forecasts, the last being the Scenario Predict-then-Optimize method of \citet{jia2025scenario}. We apply decision-focused learning at two levels. UL-DFL keeps the visit trigger of the deliver-up-to-level rule and learns how far each store's order-up-to level should deviate from the forecast quantile, depending on context and on the age of the store's stock; the learning signal is the realised holding, waste and lost-sales cost of whole training episodes, differentiated through the FIFO dynamics. LA-DFL keeps a deterministic rolling-horizon plan and learns its planning targets, differentiating the realised cost of the plan's deliveries through a coverage layer while an exact MILP oracle fixes the routes.

The answer is that learning pays when it is placed on the right structure. Across the settings studied, the stock-age-based rules with their orders routed optimally were cheaper than every integrated policy unless travel was expensive, and UL-DFL improved the best of them in every setting, significantly in eleven of thirteen, which made it the cheapest policy overall, or tied for cheapest, except with expensive routing. It does so mainly by serving the same fill rate with less waste than any parameterisation of the rule can; a rule tuned in hindsight on the test data matches its cost but not its waste. The gain is smallest where the rule already wastes little or where demand is too variable for the context to help, and it depends on the rule: the same learner without the rule's structure was never cheaper than UL-DFL and mostly not cheaper than the rule itself. Its corrections are interpretable: lower order-up-to levels on promotion days and after the weekly demand peak, when leftover stock meets falling demand. On top of the deterministic lookahead, decision-focused learning revealed that only the executed day needs a high safety quantile; a lookahead with day-specific quantile levels reproduces that gain without learning.

The paper makes three contributions.
\begin{enumerate}[leftmargin=*]
\item \textbf{Decision-focused learning of perishable order-up-to levels on shared routes.} UL-DFL learns context- and stock-age-dependent order-up-to levels for a published perishable inventory-routing rule from the realised cost of whole episodes, with model selection on the closed-loop system including exact routing. It is cheaper than the rule it starts from in every setting, significantly in eleven of thirteen, and the cheapest or joint-cheapest policy unless routing is expensive; its corrections can be read as inventory rules, and an ablation without the rule shows that the structure is what makes the learning pay. We show that the inventory recourse admits FIFO-optimal solutions once early disposal is allowed, bound the marginal value of a delivery by the lost-sales and waste penalties, and prove that stationary points of the decision-focused objective of the lookahead learner satisfy the stochastic optimality conditions of the fixed-route problem, reducing to the perishable newsvendor fractile for one-day shelf life; all results are verified numerically.
\item \textbf{An inventory-centred benchmark of replenishment policies under routing coupling.} We evaluate stock-age-based inventory rules, a published stochastic PIRP rule, and integrated deterministic, robust, scenario-based and decision-focused lookahead policies under one closed-loop simulator, with every lookahead solved by an exact MILP to a proven gap, three network layouts, 100 test episodes per setting, three training seeds, prespecified hypotheses, and a case with observed fresh-food demand. Outcomes are reported in inventory terms---cost composition, fill rate, waste rate, delivery frequency and store-level service---across shelf life, cost ratio, fleet tightness, routing cost, demand variability and issuing policy.
\item \textbf{Managerial evidence on when integration, stochastic lookahead and learning pay off.} The results show when simple stock-age-based rules suffice and when integration with routing pays, why deterministic plans over-stock and how a day-specific safety stock repairs this, which corrections of an order-up-to rule reduce waste, and how integrated plans concentrate stock-outs at individual stores (Section~\ref{sec:results}).
\end{enumerate}

Section~\ref{sec:literature} reviews related work, Section~\ref{sec:problem} defines the problem and its structural properties, Section~\ref{sec:policies} describes the policies, Section~\ref{sec:design} the experimental design, Section~\ref{sec:results} the results, and Section~\ref{sec:conclusion} concludes with managerial implications and limitations.
